\begin{abstract}
Data-management systems increasingly modify the queries users run in order to satisfy external requirements: query-refinement systems rewrite selection predicates so that the output meets a representation constraint, and consistent query answering removes the tuples involved in integrity violations.
The modified query meets the requirement but leaves the user with new questions.
Which population does the new answer describe, and how does it differ from the one the original query selected?
And when the data is sensitive, can such differences be explained without revealing information about individuals?
To address these questions, we present a framework that, to the best of our knowledge, is the first to explain the difference between two query outputs with differential privacy (DP).
An explanation is a pair of histograms of one attribute over the tuples of the two outputs that fall in a view, a conjunctive region of the data.
We adapt the measures of explanation quality used in prior work, namely representativeness, coverage, interestingness and diversity, to this setting and bound their sensitivity.
Building on them, the framework
(a)~privately selects views that are representative of the data and cover the two outputs,
(b)~privately selects from these views a diverse set of explanations whose distributions differ most between the outputs, and
(c)~releases their histograms with noise,
while guaranteeing DP in all steps.
\end{abstract}
